\documentclass{replab}
\usepackage{svg}
% --- Información del documento ---
\title{Proyecto 1}
\author{Francisco Uriel Pineda Hernández}

% Nota: Si se desea incluir más de un autor en el documento, el
% archivo replab.cls, en la sección "Página de título de documento",
% contiene líneas de código comentadas pensadas para introducir los
% datos desde 1 hasta 4 autores. Sin embargo, debe escogerse solo una
% de las cuatro secciones de código y comentar las demás para mantener
% la consistencia del documento.

\date{\today}
\subtitle={Reporte técnico de la aplicación: Chat}
\email={\href{mailto:ciscohdz@ciencias.unam.mx}{\color{principaluno}\texttt{ciscohdz@ciencias.unam.mx}}}
\subject={Modelado y programación}

\setlength{\columnsep}{14pt}

% --- Archivo de bibliografía ---
\addbibresource{repbib.bib}

% --- Inicio del documento ---
\begin{document}
	
	\pagestyle{fancy}
	\unspacedoperators
	
% --- Título ---
	\twocolumn[
		\begin{center}
			\maketitle
			
			{\begin{tcolorbox}[colframe=white, colback=principaldos, arc=8pt]
				\begin{onecolabstract}

                  Se describe el diseño e implementación de una
                  aplicación de mensajería. El proyecto se estructura
                  con la arquitectura cliente-servidor. Hace uso del
                  protocolo TCP/IP, La comunicación entre el cliente y
                  el servidor se da mediante sockets y el intercambio
                  de datos con mensajes en formato JSON. Se hace uso
                  de la programación concurrente, mediante hilos de
                  ejecución, para atender la interacción que se puede
                  dar entre el servidor y los múltiples clientes
                  conectados simultáneamente a él. El servidor,
                  escrito en el lenguaje de programación C, utiliza
                  las bibliotecas cJSON y GLib para extender las
                  funcionalidades nativas del mismo. El resultado se
                  fue la implementación de un servidor capaz de
                  soportar la carga de atender múltiples conexiones de
                  manera concurrente y servir a las peticiones
                  requeridas, siguiendo la especificación del
                  protocolo de comunicación, mientras asegura la
                  integridad de la información global de todos los
                  usuarios.

				\end{onecolabstract}
			\end{tcolorbox}}

			\medskip
		\end{center}
	]

	\selectlanguage{spanish}
	
% --- Cuerpo del reporte ---
	
	\section{Introducción}

    En los sistemas orientados a la mensajería en tiempo real, uno de
    los desafíos fundamentales radica en coordinar de forma efectiva
    el intercambio de información entre múltiples participantes de
    manera simultánea. Cuando varios usuarios intentan interactuar
    mediante mensajes privados, salas particulares o difusiones
    globales, es necesario el manejo rápido de las comuninaciones, la
    sincronización y la consistencia de los datos. Sin un mecanismo
    central de gestión, las acciones de un usuario podrían interferir
    con las de los demás o derivar en pérdidas de información.

    El objetivo principal de este proyecto es diseñar e implementar un
    sistema de chat funcional que atienda estas necesidades mediante
    la un agente central que sirva de intermediario. El alcance del
    proyecto está delimitado a un escenario de escala relativamente
    pequeña o acotada; sin embargo, el enfoque no se limita a resolver
    el flujo básico de mensajes, sino que busca garantizar un nivel
    alto de robustez. Se prioriza que la comunicación sea fiable, de
    modo que el sistema responda adecuadamente ante desconexiones
    inesperadas, mantenga la integridad de la información global y
    minimice los tiempos de espera del cliente. 
    
	\section{Modelado}

    Habiendo establecido las necesidades del sistema y los objetivos a
    cumplir, modelaremos la solución analizando detalladamente los
    requerimientos para dividir cada problema en pequeños fragmentos
    más manejables para, posteriormente, diseñar componentes que los
    resuelvan.
	
	\subsection{Análisis}

    La primera división surge de manera natural por la necesidad que
    tenemos de servir a múltiples usuarios, por lo que tomaremos
    inspiración de algunas de las aplicaciones de comunicación
    existentes, que sabemos se sustentan bajo la arquitectura
    cliente-servidor. De ahí que podamos diferenciar las
    responsabilidades que tiene uno en comparación con el otro, sin
    olvidar que dependen entre sí para lograr el objetivo de
    comunicar.

    \subsection{Analizando el Servidor}

    El servidor es, probablemente, la parte más compleja del sistema,
    ya que realiza varias acciones a la vez. De manera muy general,
    notamos que el servidor tiene que realizar las siguientes
    acciones:

    \begin{enumerate}
    \item Establecer conexiones que se mantengan activas por un tiempo
      indeterminado.
    \item Recibir peticiones, procesarlas y determinar que sean
      realizables.
    \item Realizar de manera concurrente las operaciones que le sean
      solicitadas por los clientes.
    \item Almacenar información global del sistema asegurando la
      integridad de la misma.
    \item Transmitir las respuestas y notificaciones únicamente a los
      clientes correspondientes.
    \item Exhibir un comportamiento robusto ante fallos imprevistos o
      errores en la comunicación.
    \end{enumerate}

    Del primer punto, notamos que el mantener la conexión es algo que
    no depende enteramente del servidor. Bien puede pasar que el servidor
    finalice su ejecución, así como que un usuario decida desconectarse o,
    en casos más extremos, que la conexión entre ambos falle por motivos
    que escapan de nuestro control. En todos los casos anteriores, el
    servidor debe poder reaccionar para eliminar el registro de ese
    cliente y dar los avisos correspondientes.

    Las acciones que el servidor puede realizar están especificadas en el
    protocolo de comunicación; es por eso que el punto dos está bien
    delimitado. Es primordial hacer que dicho protocolo se cumpla tanto al
    recibir como al enviar peticiones/notificaciones, ya que de esa manera
    aseguramos la correcta comunicación entre un cliente/servidor ajeno a
    este proyecto, pero que sigue la especificación del protocolo.

    Atender las solicitudes de manera concurrente quiere decir que, en
    lugar de realizar cada acción de inicio a fin una detrás de la
    otra, hacemos pequeños pedazos de cada tarea para avanzar en
    todas. Veremos que al usar concurrencia reducimos el tiempo de
    espera que un cliente sufre al enviar/recibir mensajes, pero
    mencionaremos más adelante posibles problemas que surgen de esta
    estrategia.

    Es de suma importancia que el servidor tenga el registro de los
    clientes que debe comunicar, de las salas a las que pertenecen y
    del modo en que puede hacerles llegar notificaciones. De la misma
    forma, debe asegurar que los datos registrados sean correctos, es
    decir, que correspondan a los clientes correctos y que reflejen
    las acciones que han realizado a lo largo del tiempo. Tenemos un
    problema al aplicar concurrencia ya que puede ocurrir que la
    información sea modificada de manera incorrecta, esto ocurre
    debido a que la división solo aplica a las acciones a realizar y
    no a lo que se almacena, por lo que, si en una tarea se cambia o
    elimina algún dato, es casi seguro que al intentar leer ese dato,
    sin saber de su modificación, terminemos ocasionando errores
    graves en la comunicación.

    Revisamos de manera indirecta los puntos cinco y seis en los
    párrafos anteriores, en los que hablamos de la transmisión y
    recepción, así como de la precauciones a tener en cuenta con
    ciertos errores que podemos intuir se darán.

    Las acciones específicas que determina el protocolo se detallarán
    en el diseño, donde podremos determinar la manera de manejarlas y
    encapsularlas junto a las necesidades aquí descritas.

    \subsection{Analizando el Cliente}

    El cliente, al ser una pieza controlada enteramente por un usuario,
    nos conviene dividirlo en tres:

    \begin{enumerate}
    \item La parte más cercana al usuario, que simplificará la
      comunicación con el servidor, haciendo uso de comandos más sencillos
      de escribir y recordar.
    \item La parte más cercana al servidor, que establecerá conexión y
      sabrá cómo comunicarse con él.
    \item Un intermediario entre los dos anteriores, que sepa traducir
      los comandos del usuario a los mensajes que se envían al
      servidor y viceversa.
    \end{enumerate}
    
    El primer componente, al ser una pieza enfocada a la comodidad del
    usuario, debe tener un lenguaje de comandos simples que ayuden a
    determinar rápidamente la operación a realizar. Además, el cómo
    presentamos las diferentes notificaciones que el servidor envía
    debe ser claro para quienes estén dirigidas. Exploraremos en el
    diseño la manera de hacer que esto ocurra.

    El cliente debe poder conectarse a un servidor, así que se debe
    notar la necesidad de establecer y mantener una conexión activa
    con este último. De manera casi opuesta a lo que verá el usuario,
    para enviar y recibir mensajes al servidor, la segunda parte debe
    hacer uso enteramente del protocolo de comunicación, por lo que
    debe saber armar los mensajes con el formato específico y recibir
    las notificaciones del servidor.

    Para completar el viaje de una petición de usuario, necesitamos
    traducir el comando que llega desde la primera parte hasta la
    segunda. Por tanto, la tercera parte del cliente será un traductor
    entre las dos primeras. Notamos que es probablemente la más compleja,
    ya que tiene que determinar: la correcta escritura de comandos; enviar
    los datos necesarios para que la parte más cercana al servidor forme
    los mensajes a enviar; traducir la notificación que la segunda parte
    recibió y enviarla en un formato presentable a la primera.

    Un detalle importante a notar es que necesitamos enviar peticiones
    al servidor mientras, se reciben notificaciones y se le presentan
    al usuario. Para ello usaremos concurrencia. Al igual que en el
    servidor, la ejecución concurrente debe gestionarse adecuadamente
    para evitar colisiones entre el flujo de lectura de notificaciones
    y la entrada de comandos del usuario.

    \begin{figure}[htbp]
		\centering
        \includesvg[width=0.5\textwidth]{img/cliente_servidor.svg}
		\caption{Diagrama de la división cliente servidor.}
		\label{fig:cliente_servidor}
	\end{figure}

    \subsection{Diseño}

    Teniendo la división que vemos en la Figura
    \ref{fig:cliente_servidor}, nos será mas senciilo identificar los
    componentes del sietema y de qué forma deberán comunicarse.
    
    \subsection{Diseño del servidor}

    Podemos agrupar las necesidades del servidor, previamente
    analizadas, en 'capas'. La Figura \ref{fig:div_serv} ofrece una
    representación visual de dicha división.

    \begin{figure}[htbp]
		\centering
        \includesvg[width=0.5\textwidth]{img/div_serv.svg}
		\caption{Diagrama de la división del servidor.}
		\label{fig:div_serv}
	\end{figure}

    \subsubsection{Capa de Comunicación}

    Se encarga de todo lo relativo a la conexión y comunicación con
    clientes asi como a la interpretación y estructuración de mensajes
    que siguen la especificación del protocolo.

    Notamos entonces que podemos dividir la capa en dos componentes.

    El primero siendo aquel que sabe establecer una conexión con un
    cliente. La conexión consta, a su vez, de dos partes: se debe
    identficar a un cliente para comunicarse con él y se debe poder
    enviar/recibir información por medio de la conexión. 

    Como no podemos recibir cualquier mensaje, el segundo componente
    deberá ser capaz de determinar si un mensaje está bien formado,
    esto es, que contenga el formato especificado en el protocolo
    aunque, como la capa de comunicación no se comunica directamente
    con la de almacenamiento, es imposible determinar si un mensaje es
    realizable pero sí podemos determinar qué tipo de petición estamos
    recibiendo. Al seguir en la capa más próxima a la interacción con
    el cliente, nos conviene que el componente también sea capaz de
    estructurar los mensajes de notificación.

    \subsubsection{Capa de Operaciones}
    
    Es la responsable de dirigir el flujo de trabajo del servidor,
    determinando las operaciones que se realizan y la modificación de
    los registros de los clientes. Además debe manejar, junto con la
    \textit{capa de comunicación}, el envío de mensajes.

    Sobre las operaciones a realizar, notamos que podemos tener un
    punto de entrada sobre el cual, una petición verificada por la
    \textit{capa de comunicación}, se distribuye dependiendo el tipo
    de la misma o se determina como inválida dentro de los límites del
    protocolo.

    Como se mencionaba en el análisis, todas las peticiones necesitan,
    en mayor o menor medida, de la información de la \textit{capa de
      almacenamiento} para poder realizarse.

    Dado que tenemos el problema de varias tareas realizándose a la
    vez, podemos tener problemas inesperedados de modificación y
    lectura sobre los datos que se comparten entre las distintas
    tareas. Una forma de solucionar dicho problema es cada que una
    tarea quiere acceder o modificar la información compartida tenga
    una manera de saber si alguna otra lo está haciendo en primer
    lugar y entonces esperar a que el recurso sea `liberado`. A esta
    técnica se le conoce como \textit{exclusión mutua} y continuaremos
    discutiendo las partes que protegeremos de los fallos de la
    concurrencia en la capa de almacenamiento y en la implementación
    del diseño. 
    
    \subsubsection{Capa de Almacenamiento}
    
    Por último en esta capa se encuentran toda las estructuras que
    guardan la información relativa a los clientes. Aquí se establece
    la manera es que se crean y eliminan dichas estructuras.

    Es primordial que guardemos a cada cliente que haya establecido
    conexión con el servidor y se haya identificado apropiadamente.

    Un primer acercamiento a esta necesidad nos deja en claro que cada
    cliente tiene datos que necesitamos mantener accesibles: la
    identificación de un cliente para poder enviar información
    mediante la conexión que establecimos con él; un identificador
    conocido por los demás clientes conectados; y, siguiendo el
    protocolo, su estado en el servidor (ACTIVE, BUSY, AWAY).

    Se tienen que almacenar los clientes en una estructura que nos
    permita agregar a cuántos usuarios decidan conectarse al servidor
    y eliminarlos cuando se desconecten. Además, en el protocolo se
    especifica la posibilidad de crear salas privadas, dónde solo
    algunos usuarios, del total de conectados, pueden enviar mensajes
    que solo recibirán ellos. Además pueden existir múltiples salas en
    el servidor. Por ello también se contempla la creación de una
    estructura que guarde salas (dónde sobra decir que pueden agregar
    nuevas salas y eliminarlas cuando dejen de tener usuarios) dónde
    cada sala es, a su vez, una estructura que guarda usuarios dentro
    de la y aquellos invitados a estar dentro de ella.

    Visto de forma global, podemos crear una estructura que almacena a
    los clientes de todo el servidor y las salas que en él se crearon
    cada uno con sus respectivos clientes.

    Sin dejar a un lado que debemos proteger a toda esta estructura
    global de las múltiples modificaciones que pueden resultar de la
    concurrencia. Discutiremos de forma concreta qué hacemos para
    lograr este objetivo en los detalles de la implementación. 

	\section{Implementación}

    Se opta por el lenguaje de programación C para implementar el
    servidor ya que este, a pesar de las pocas facilidades que ofrece
    al programador, resulta sumamente enriquecedor para entender
    conceptos fundamentales como el manejo de memoria dinamica.

    Sin embargo el objetivo no es reinventar la rueda. Por este motivo
    también se usan tecnologías y bibliotecas de terceros que ayudan
    enormemente al desarrollo de un proyecto de esta magnitud.

    Empezando por el sistema de contrucción que ayuda a mantener
    centralizado todo lo relativo a la compilación, manejo de
    dependencias y organización general del proyecto. En el caso de la
    implementación presentada se usa el sistama de contrucción:
    \textit{Meson}, el cual maneja todo lo necesario para la
    instalación de las dependencias necesarias asi como de la
    compilación del proyecto entero. Resulta bastante intuitivo una
    vez que se explora con cierto detalle las funcionalidades básicas
    que posee.

    Hoy en un día, usar un manejador de versiones para llevar el
    progreso de cualquier proyecto, es un estándar. Con diferencia el
    más popular es \textbf{git} y es el que se usa para llevar un
    control de los cambios y modificaciones en el desarrollo del
    sistema.

    Las bibliotecas usadas para el servidor son \textbf{GLib} y
    \textbf{cJSON}. La primera implementa y extiende muchas de la
    funcionalidades de la biblioteca estándar de C, en particular se
    usa la implemetanción de \textit{GLib} de diccionarios
    (GHashTable) para cumplir con el diseño de la \textit{Capa de
      Almacenamiento} del servidor que se modela en la Figura
    \ref{fig:estruct_serv}

    \begin{figure}[htbp]
		\centering
        \includesvg[width=0.5\textwidth]{img/estruct_serv.svg}
		\caption{Diagrama de las estructuras del servidor siguiendo el diseño propuesto.}
		\label{fig:estruct_serv}
	\end{figure}

    El diseño concurrente se implementa usando los hilos de ejecución
    propios de la biblioteca \textit{GLib}, los mismos se encunetran
    en el punto de entrada del servidor y se crea uno nuevo al
    establecer la conexión con un cliente. Para mantener la consistencia de
    la información almacenada se implementó la exclusión mutua en la
    estructura \textit{ContextoServidor}, haciendo uso de la
    estructura \textit{GMutex}.

    Para establecer conexiones y cumplir con el manejo correcto del
    protocolo de comunicación, siguiendo el diseño de la \textit{Capa
      de Comunicación}, se usa la API de POSIX para sockets en C, que
    proveen de todo lo necesario para establecer conexiones en redes
    locales y comunicarse haciendo uso del protocolo TCP/IP, ideal
    para el tipo de sistema que realizamos.

    Un detalle de implemanteación es que debido a lo largo y
    complicado que puede ser el codigo necesario para llevar a cabo
    las conexiones haciendo uso de sockets se optó por envolverla en
    un solo componente, \textit{$socket_util$}, que establece las
    funciones para inciar conexiones, esperar un conexión y enviar
    información a traves de una.
    
    Se crea además un componente, \textit{protocolo}, que define
    enumeraciones y mapeos a cadenas para poder identificar
    rápidamente los tipos de mensajes, operaciones, resultados y
    estados que se pueden enviar en un JSON. A su vez, usando los
    patrones de diseño \textbf{DTO} (Data Transfer Object) y
    \textit{Factory} se crea una estructura \textit{Campos} que se
    pasa como una \textbf{Literal Compuesta} (Compound Literals) en C
    a la fución \textit{$fabrica_respuesta$} que usa la biblioteca
    \textit{cJSON} para traducir y empaquetar la información necesaria
    en la comunicación.

    En la \textit{Capa de Operaciones}, tenemos a los componentes
    \textit{distribuidor} y \textit{manejador} que se encargar,
    dependiendo el tipo de mensaje, de realizar las peticiones
    necesarias y enviar las notificaciones correspondientes en cada
    caso.

    \section{Conclusión}

    Se logró abstraer con éxito la complejidad del desarrollo de un
    servidro en C mediante el uso de tipos abstractos de datos,
    patrones de diseño y la biblioteca GLib, alcanzando un servidor de
    chat multihilo estable, capaz de gestionar múltiples salas y
    usuarios de forma concurrente.

    Fue un gran proyecto que ante todo marca un antes y un después en
    lo que creí que me tomaría años aprender. A la vez resulta ser uno
    de los que más satisfaccción me deja a pesar de las complicaciones
    y contratiempos que hicieron que no pudiera terminarlo como me
    gustaría. 

    Fue abrumador balancear el análisis y el diseño de cosas que no
    terminaba de entender por qué eran necesarias o cómo
    funcionaban. Y, una vez puesto en práctica, los erores referentes
    a las bibliotecas desconocidas resultaban abrumadores.

    Me llevo mucha experiencia y ganas de continuar haciendo proyecto
    que no me veía haciendo antes de entrar a la Licenciatura.

    \printbibliography[heading=bibintoc]
	
\end{document}
